% GENERATED FILE. Edit canonical lessons, then run npm run generate:books.
% canonical-source-hash: fnv1a64:4c551ff6b5f04b7f
% canonical-lessons: PA-C352-shatranj, PA-C352-bazi, PA-C352-vari, PA-C352-goti, PA-C352-inam

\chapter{Games Indoors}
\label{ch:pa-a2-games-indoors}
\hlchaptermodality{352}

\begin{chapteropening}
\textbf{By the end of this chapter:} \emph{I can talk about the indoor games I like.}

Complete the last lesson of chapter 352: I can talk about the indoor games I like.
\end{chapteropening}

\section[\texorpdfstring{\pa{ਸ਼ਤਰੰਜ} (shatranj)}{ਸ਼ਤਰੰਜ (shatranj)}]{\pa{ਸ਼ਤਰੰਜ} (shatranj) — chess}

\label{lesson:PA-C352-shatranj}



\begin{quote}
Before the new one: say the Punjabi for a match, then the Punjabi for to defeat.
\end{quote}

\subsection*{You'll want to know: \pa{ਸ਼ਤਰੰਜ}}
\textbf{\pa{ਸ਼ਤਰੰਜ}} — \emph{shatranj} — \textquotedblleft{}chess\textquotedblright{}.

\begin{quote}
\pa{ਮੇਰਾ} \pa{ਦਾਦਾ} \pa{ਸ਼ਾਮ} \pa{ਨੂੰ} \pa{ਸ਼ਤਰੰਜ} \pa{ਖੇਡਦਾ} \pa{ਹੈ।}
\end{quote}

\begin{quote}
\emph{merā dādā shām nū̃ shatranj kheḍdā hai.}
\end{quote}

\begin{quote}
My grandfather plays chess in the evening.
\end{quote}

A feminine noun.

\begin{cousinweb}
\textbf{\pa{ਸ਼ਤਰੰਜ}} (\emph{shatranj}) is a loan, usually traced to Persian \emph{shatranj}, \textquotedblleft{}chess\textquotedblright{}, which goes back to Sanskrit \emph{caturaṅga}, \textquotedblleft{}four-limbed\textquotedblright{}.
\end{cousinweb}

\begin{grammarlens}[title={What you've built}]
One more word for this chapter.
\end{grammarlens}

\subsection*{Your turn}
\begin{itemize}
\raggedright
  \item \emph{Say it:} \emph{shatranj}
  \item \emph{Say it:} \emph{shatranj}, once more
  \item \emph{Say it:} \emph{merā dādā shām nū̃ shatranj kheḍdā hai.} — My grandfather plays chess in the evening.
  \item \emph{Recall:} say the Punjabi for a match, then the Punjabi for to defeat, then say \emph{shatranj} again
\end{itemize}

\begin{tcolorbox}[breakable,colback=teal!4,colframe=teal!35!black,title={Before you move on}]
What does \pa{ਸ਼ਤਰੰਜ} mean? (\textquotedblleft{}Chess\textquotedblright{}.) Say it once more.
\end{tcolorbox}
\section[\texorpdfstring{\pa{ਬਾਜ਼ੀ} (bāzī)}{ਬਾਜ਼ੀ (bāzī)}]{\pa{ਬਾਜ਼ੀ} (bāzī) — a round (of a game); a game}

\label{lesson:PA-C352-bazi}



\begin{quote}
Before the new one: say the Punjabi for to defeat, then the Punjabi for chess.
\end{quote}

\subsection*{You'll want to know: \pa{ਬਾਜ਼ੀ}}
\textbf{\pa{ਬਾਜ਼ੀ}} — \emph{bāzī} — \textquotedblleft{}a round (of a game); a game\textquotedblright{}.

\begin{quote}
\pa{ਅਮਨ} \pa{ਨੇ} \pa{ਪਹਿਲੀ} \pa{ਬਾਜ਼ੀ} \pa{ਜਿੱਤੀ।}
\end{quote}

\begin{quote}
\emph{Aman ne pahilī bāzī jittī.}
\end{quote}

\begin{quote}
Aman won the first round.
\end{quote}

A feminine noun.

\begin{cousinweb}
\textbf{\pa{ਬਾਜ਼ੀ}} (\emph{bāzī}) is a loan, usually traced to Persian \emph{bāzī}, \textquotedblleft{}play\textquotedblright{}.
\end{cousinweb}

\begin{grammarlens}[title={What you've built}]
One more word for this chapter.
\end{grammarlens}

\subsection*{Your turn}
\begin{itemize}
\raggedright
  \item \emph{Say it:} \emph{bāzī}
  \item \emph{Say it:} \emph{bāzī}, once more
  \item \emph{Say it:} \emph{Aman ne pahilī bāzī jittī.} — Aman won the first round.
  \item \emph{Recall:} say the Punjabi for to defeat, then the Punjabi for chess, then say \emph{bāzī} again
  \item \emph{Write it:} copy \textbf{\pa{ਬਾਜ਼ੀ}} once, checking the whole word
\end{itemize}

\begin{tcolorbox}[breakable,colback=teal!4,colframe=teal!35!black,title={Before you move on}]
What does \pa{ਬਾਜ਼ੀ} mean? (\textquotedblleft{}A round (of a game); a game\textquotedblright{}.) Say it once more.
\end{tcolorbox}
\section[\texorpdfstring{\pa{ਵਾਰੀ} (vārī)}{ਵਾਰੀ (vārī)}]{\pa{ਵਾਰੀ} (vārī) — one's turn (in a game or a queue)}

\label{lesson:PA-C352-vari}



\begin{quote}
Before the new one: say the Punjabi for chess, then the Punjabi for a round of a game.
\end{quote}

\subsection*{You'll want to know: \pa{ਵਾਰੀ}}
\textbf{\pa{ਵਾਰੀ}} — \emph{vārī} — \textquotedblleft{}one's turn (in a game or a queue)\textquotedblright{}.

\begin{quote}
\pa{ਹੁਣ} \pa{ਮੇਰੀ} \pa{ਵਾਰੀ} \pa{ਹੈ।}
\end{quote}

\begin{quote}
\emph{huṇ merī vārī hai.}
\end{quote}

\begin{quote}
Now it's my turn.
\end{quote}

A feminine noun.

\begin{cousinweb}
\textbf{\pa{ਵਾਰੀ}} (\emph{vārī}) is inherited from Sanskrit \emph{vāra}, \textquotedblleft{}appointed time, one's turn\textquotedblright{}; CDIAL lists Punjabi \emph{vār}, \textquotedblleft{}occasion\textquotedblright{}, and Lahnda \emph{vārī}, \textquotedblleft{}turn, time\textquotedblright{}, under it (CDIAL 11547).
\end{cousinweb}

\begin{grammarlens}[title={What you've built}]
One more word for this chapter.
\end{grammarlens}

\subsection*{Your turn}
\begin{itemize}
\raggedright
  \item \emph{Say it:} \emph{vārī}
  \item \emph{Say it:} \emph{vārī}, once more
  \item \emph{Say it:} \emph{huṇ merī vārī hai.} — Now it's my turn.
  \item \emph{Recall:} say the Punjabi for chess, then the Punjabi for a round of a game, then say \emph{vārī} again
  \item \emph{Write it:} copy \textbf{\pa{ਵਾਰੀ}} once, checking the whole word
\end{itemize}

\begin{tcolorbox}[breakable,colback=teal!4,colframe=teal!35!black,title={Before you move on}]
What does \pa{ਵਾਰੀ} mean? (\textquotedblleft{}One's turn (in a game or a queue)\textquotedblright{}.) Say it once more.
\end{tcolorbox}
\section[\texorpdfstring{\pa{ਗੋਟੀ} (goṭī)}{ਗੋਟੀ (goṭī)}]{\pa{ਗੋਟੀ} (goṭī) — a game piece; a counter on a board}

\label{lesson:PA-C352-goti}



\begin{quote}
Before the new one: say the Punjabi for a round of a game, then the Punjabi for one's turn.
\end{quote}

\subsection*{You'll want to know: \pa{ਗੋਟੀ}}
\textbf{\pa{ਗੋਟੀ}} — \emph{goṭī} — \textquotedblleft{}a game piece; a counter on a board\textquotedblright{}.

\begin{quote}
\pa{ਮੇਰੇ} \pa{ਕੋਲ} \pa{ਚਾਰ} \pa{ਲਾਲ} \pa{ਗੋਟੀਆਂ} \pa{ਹਨ।}
\end{quote}

\begin{quote}
\emph{mere kol chār lāl goṭīāṁ han.}
\end{quote}

\begin{quote}
I have four red pieces.
\end{quote}

A feminine noun.

\begin{cousinweb}
\textbf{\pa{ਗੋਟੀ}} (\emph{goṭī}) is inherited from a reconstructed \emph{*goṭṭa}, \textquotedblleft{}something round\textquotedblright{}; CDIAL lists Punjabi \emph{goṭ}, \textquotedblleft{}piece on a chequer board\textquotedblright{}, under it (CDIAL 4271), but not the form \emph{goṭī} itself.
\end{cousinweb}

\begin{grammarlens}[title={What you've built}]
One more word for this chapter.
\end{grammarlens}

\subsection*{Your turn}
\begin{itemize}
\raggedright
  \item \emph{Say it:} \emph{goṭī}
  \item \emph{Say it:} \emph{goṭī}, once more
  \item \emph{Say it:} \emph{mere kol chār lāl goṭīāṁ han.} — I have four red pieces.
  \item \emph{Recall:} say the Punjabi for a round of a game, then the Punjabi for one's turn, then say \emph{goṭī} again
  \item \emph{Write it:} copy \textbf{\pa{ਗੋਟੀ}} once, checking the whole word
\end{itemize}

\begin{tcolorbox}[breakable,colback=teal!4,colframe=teal!35!black,title={Before you move on}]
What does \pa{ਗੋਟੀ} mean? (\textquotedblleft{}A game piece; a counter on a board\textquotedblright{}.) Say it once more.
\end{tcolorbox}
\section[\texorpdfstring{\pa{ਇਨਾਮ} (inām)}{ਇਨਾਮ (inām)}]{\pa{ਇਨਾਮ} (inām) — a prize; a reward}

\label{lesson:PA-C352-inam}



\begin{quote}
Before the new one: say the Punjabi for one's turn, then the Punjabi for a game piece.
\end{quote}

\subsection*{You'll want to know: \pa{ਇਨਾਮ}}
\textbf{\pa{ਇਨਾਮ}} — \emph{inām} — \textquotedblleft{}a prize; a reward\textquotedblright{}.

\begin{quote}
\pa{ਸਕੂਲ} \pa{ਹਰ} \pa{ਸਾਲ} \pa{ਚੰਗੇ} \pa{ਵਿਦਿਆਰਥੀਆਂ} \pa{ਨੂੰ} \pa{ਇਨਾਮ} \pa{ਦਿੰਦਾ} \pa{ਹੈ।}
\end{quote}

\begin{quote}
\emph{sakūl har sāl chaṅge vidiārthīāṁ nū̃ inām dindā hai.}
\end{quote}

\begin{quote}
Every year the school gives prizes to good students.
\end{quote}

A masculine noun.

\begin{cousinweb}
\textbf{\pa{ਇਨਾਮ}} (\emph{inām}) is a loan, usually traced through Persian to Arabic \emph{inʿām}, \textquotedblleft{}favour\textquotedblright{}.
\end{cousinweb}

\begin{grammarlens}[title={What you've built}]
One more word for this chapter.
\end{grammarlens}

\subsection*{Your turn}
\begin{itemize}
\raggedright
  \item \emph{Say it:} \emph{inām}
  \item \emph{Say it:} \emph{inām}, once more
  \item \emph{Say it:} \emph{sakūl har sāl chaṅge vidiārthīāṁ nū̃ inām dindā hai.} — Every year the school gives prizes to good students.
  \item \emph{Recall:} say the Punjabi for one's turn, then the Punjabi for a game piece, then say \emph{inām} again
  \item \emph{Write it:} copy \textbf{\pa{ਇਨਾਮ}} once, checking the whole word
\end{itemize}

\begin{tcolorbox}[breakable,colback=teal!4,colframe=teal!35!black,title={Before you move on}]
What does \pa{ਇਨਾਮ} mean? (\textquotedblleft{}A prize; a reward\textquotedblright{}.) Say it once more.
\end{tcolorbox}
